\section{Conclusion}

We completed a reproducible few-shot anomaly-detection study on six VisA
categories. The template baseline demonstrates why raw pixel differences are
unreliable under pose and appearance variation. AnomalyDINO provides a strong
training-free detector by matching DINOv2 patch features against a small normal
memory. Our multi-layer variant tests whether combining semantic and spatial
transformer blocks improves the detection-localization trade-off while keeping
all other factors fixed.

The project also provides an auditable experimental protocol: the 20
calibration images are excluded from reference selection, exact filenames are
saved, test labels are used only for final scoring, and every method shares the
same references and seeds. This design supports both the quantitative report
and the interactive OK/NG demonstration.

Future work should batch the frozen encoder for higher throughput, learn a
small calibration-only weighting over transformer layers, and add selective
prediction so that ambiguous products are deferred to human inspection instead
of being forced into OK or NG.
